Avalikud äratussõna raamistikud nagu \texttt{openWakeWord} \cite{openwakeword2026} ja Picovoice Porcupine \cite{picovoice-benchmark2026} ei paku eelnevalt treenitud eestikeelseid mudeleid, mistõttu eestikeelsel kasutajal puudub privaatne lokaalne nutikodu hääljuhtimise satelliit, ehkki eesti kõnetuvastuse korpused \cite{taltech-asr-data2024} ja rahvusvahelised äratussõna võrdluspunktid \cite{microwakeword2026} on olemas. Äratussõna tuvastus on kogu hääljuhtimise esimene filter: selle vead --- nii puuduvad tuvastused kui valesvallandumised --- kanduvad muutmata edasi STT-le, dialoogiloogikale ja seadme käitumisele.

Käesoleva töö eesmärk on uurida, kuidas ehitada eestikeelse äratussõna tuvastamist piiratud ressursiga mikrokontrolleril nii, et lahendus oleks lokaalne, reprodutseeritav ja reaalses kasutuses hinnatav. Praktiliseks sihtplatvormiks on ESP32-S3-põhine seade, kuhu mudel paigaldatakse ESPHome'i tarkvarakihi kaudu \cite{esphome2026}; mudeli treenimiseks ja hindamiseks kasutatakse \texttt{microWakeWord}i raamistikku \cite{microwakeword2026,tensorflow2015}. Nutikodu poolel seotakse prototüüp Home Assistanti \cite{homeassistant2026} \texttt{voice\_assistant} liidesega.

Selle tööriistaahela rakendamise käigus ilmnes kiiresti, et väikese keele äratussõna projekti puhul ei piisa ainult mudeli treenimisest: piiratud andmestik muudab hindamisprotokolli ja tõendamise rangus keskseks. Sama mudel võib anda lühikestel testklippidel häid tulemusi, kuid pidevas helivoos valesti vallanduda. Seetõttu tuleb klassikaliste mõõdikute kõrval mõõta valevallandumisi pika taustahelisalvestisega mõõdikuga FAPH (valevallandumised tunnis) \cite{lopezespejo2021deepkws}. Sihiks seati FAPH~$<$~1. See on leebem kui küpsemate raamistike avaldatud sihid (openWakeWord FAPH~$<\!0{,}5$, Picovoice FAPH~$\approx$~0{,}1). Valik tugineb kahel kaalutlusel: (i) praeguse eestikeelse taustakorpuse pikkusega saab usaldusväärselt eristada üks vallandumine tunnis suurusjärgu sagedusi, samas kui väiksemate sageduste statistiliseks tõendamiseks oleks vaja oluliselt mahukamat korpust; (ii) töö positsioneerib end eestikeelse äratussõna \emph{metoodika alusena} --- pidevvoo hindamisahela reprodutseeritav rakendamine on olulisem kui rahvusvahelise tipptulemuse saavutamine.

Töö põhiküsimus on järgmine: kuidas luua ja hinnata eestikeelse äratussõna tuvastamist nii, et see vastaks samaaegselt madala valevallandumismäära (FAPH~$<$~1), piisava lähikõne tuvastamismäära ($\geq$~0{,}95) ja mikrokontrolleri ressursipiirangute nõuetele? Selle küsimuse lahendamiseks tuleb vastata mitmele alamküsimusele:
\begin{itemize}
    \item kuidas valideerida treeningu- ja hindamisahel enne eestikeelse andmestiku juurde liikumist;
    \item millist rolli mängivad positiivsed, negatiivsed ja taustaheliandmed äratussõna mudeli kvaliteedi hindamisel;
    \item kuidas eristada andmestikust tulenevaid probleeme tööahela tehnilistest piirangutest;
    \item kas treenitud mudel saavutab eestikeelsel taustaheli korpusel pidevvoo FAPH~$<$~1 sihi;
    \item kuivõrd ulatub lähikõne tuvastamismäär $\geq$~0{,}95 sihini ning millised on selle peamised piirangud.
\end{itemize}

Töö peamine panus on FAPH-põhine pidevvoo hindamisprotokoll ning sellega seotud reprodutseeritav eestikeelse äratussõna treeningu- ja hindamisahel. Seda illustreerivad prototüüpmudel \enquote{Kuule Kratt} ning ESP32-S3 ja Home Assistanti integratsioon lokaalse nutikodu satelliidi näitel.

\paragraph{Empiirilise ulatuse piirang.}\label{sec:scope-and-limitations}
Käesoleva töö empiiriline ulatus tuleb lugeda kui \emph{prototüübi ja hindamismetoodika} dokumentatsiooni, mitte kui empiirilist hinnangut sellele, kuidas süsteem käitub laia kasutajaskonna käes. Kõik üldistused reaalsete kõnelejate käitumisele põhinevad kahel autori-lähedasel kõnelejal (autor, Kõneleja B) ning ühel treeningust eraldi hoitud kõnelejal (Kõneleja D, $N=145$ klippi); kavandatud kuni 10 osalejaga kasutajatesti piisavat lõppandmestikku selle töö raames ei moodustatud (vt §\ref{sec:user-test-results}). Töö ei käsitle täielikku STT/TTS-ahelat ning hindamine piirdub ühe äratusfraasiga.

Töö on üles ehitatud järgmiselt. Metoodika peatükk kirjeldab kasutatud andmekorpusi, tarkvarakomponente ja hindamisloogikat. Tulemuste peatükk esitab treeningu- ja hindamisahela valideerimise ning eestikeelsete mudelikatsete tulemused. Arutelu ja järelduste peatükk seob leiud praktiliste riskide, piirangute ja lõppjäreldustega. Kokkuvõte koondab töö põhivastused.
